% Chapter 1 of The Replica Rosetta Stone: free entropy, saddle points and self-averaging.
\newcommand{\freeentropymmse}{\mathrm{mmse}}
\newcommand{\freeentropyEnt}{\operatorname{Ent}}

\rstchapter{Free Entropy}{Posteriors, saddle points and self-averaging}{saddle points\dict large deviations}{ch:freeentropy}
\stonefoot{proved}
\begin{refsection}[chapters/01-freeentropy.bib]

\begin{chaptersummary}
A Bayesian posterior over $N$ variables is a Gibbs measure, its normalizer is a partition function, and the expected log-evidence per variable is a free entropy, minus the mutual information per variable up to an explicit term. Physicists compute free entropies at a saddle point, average over the data with replicas, and expect the answer to self-average. We set this beside the rigorous route on three models: the Curie--Weiss model, where the saddle point is Laplace's method with error at most $\log(N+1)/N$; the Sherrington--Kirkpatrick model, where self-averaging is Gaussian concentration and integer replica moments are a large-deviation computation, but the continuation to zero replicas has no counterpart; and the scalar Gaussian channel, where Bayes' rule gives the Nishimori identity and integration by parts the I-MMSE relation. Exact enumeration and quadrature confirm every bound and expose the rare samples that dominate integer moments; the size of the fluctuations at low temperature is open.
\end{chaptersummary}

\section{The question: the evidence of a posterior}\label{sec:freeentropy:question}

A statistician observes a noisy symmetric matrix and wants the signs $x^*\in\{\pm1\}^N$ that produced it,
\begin{equation}\label{eq:freeentropy:z2}
A=\frac\lambda Nx^*x^{*\top}+W,\qquad W\sim\mathrm{GOE}(N),
\end{equation}
with $x^*$ uniform and a signal strength $\lambda\ge0$: $\Z_2$ synchronization, the spiked Wigner model with a Rademacher prior (\cref{ch:nishimori,ch:spiked}). Since the $W_{ij}$, $i<j$, are independent $\Normal(0,1/N)$, expanding the squares in the likelihood gives the posterior
\begin{equation}\label{eq:freeentropy:z2post}
\mu_A(x)=\frac1{Z_N(A)}\exp\Bigl(\lambda\sum_{i<j}A_{ij}x_ix_j\Bigr),
\end{equation}
with $Z_N(A)$ the sum of the exponential over $x\in\{\pm1\}^N$; the diagonal of $A$ carries no information, since $x_i^2=1$.

A physicist reads \eqref{eq:freeentropy:z2post} as a Gibbs measure of $N$ Ising spins with random couplings. Write $A_{ij}=\frac\lambda Nx^*_ix^*_j+g_{ij}/\sqrt N$ with independent standard Gaussian $g_{ij}$, and compare each candidate with the truth through $\tau_i=x_ix^*_i$. The $\tilde g_{ij}=g_{ij}x^*_ix^*_j$ are again independent standard Gaussians, independent of $x^*$, and $\sum_{i<j}\tau_i\tau_j=\frac12((\sum_i\tau_i)^2-N)$, so
\begin{equation}\label{eq:freeentropy:split}
\lambda\sum_{i<j}A_{ij}x_ix_j=\underbrace{\frac\lambda{\sqrt N}\sum_{i<j}\tilde g_{ij}\tau_i\tau_j}_{\text{noise}}+\underbrace{\frac{\lambda^2}{2N}\Bigl(\sum_i\tau_i\Bigr)^2}_{\text{signal}}-\frac{\lambda^2}2 .
\end{equation}
The noise is the Hamiltonian of the Sherrington--Kirkpatrick model (SK, \cref{ch:sk}) at inverse temperature $\beta=\lambda$; the signal is that of the Curie--Weiss model, the simplest mean-field model of a ferromagnet, at $\beta=\lambda^2$, and its magnetization $\frac1N\sum_i\tau_i$ is the overlap of $x$ with the truth. The two models of this chapter are the two halves of one posterior.\sidenote{\textcite{SherringtonKirkpatrick1975} took couplings of mean $J_0/N$ and variance $J^2/N$; \eqref{eq:freeentropy:split} is their model on the line $\beta J_0=(\beta J)^2$, the Nishimori line of \cref{sec:freeentropy:replicas}.}

\begin{definition}[Posterior, Gibbs measure, free entropy]\label[definition]{def:freeentropy:gibbs}
Let $P_0$ be a prior on a set $\cX$, and let the data $Y$ have the law $P(\cdot\mid x)$ when the unknown is $x\in\cX^N$, with a density with respect to a reference law $Q$ free of $x$. The \term{Hamiltonian}\index{Hamiltonian} is the log-likelihood ratio $H_N(x;Y)=\log\frac{\dd P(\cdot\mid x)}{\dd Q}(Y)$; the posterior is the \term{Gibbs measure}\index{Gibbs measure} $\mu_Y(\dd x)=e^{H_N(x;Y)}P_0^{\otimes N}(\dd x)/Z_N(Y)$, and the \term{partition function}\index{partition function} $Z_N(Y)=\int e^{H_N(x;Y)}P_0^{\otimes N}(\dd x)$ is the evidence relative to $Q$. With $X\sim P_0^{\otimes N}$ and $Y\sim P(\cdot\mid X)$, the \term{free entropy}\index{free entropy} is the expected log-evidence per variable $\phi_N=\frac1N\E\log Z_N(Y)$, with limit $\phi$ when it exists. The \term{Gibbs average} of $f$ is $\gibbs f=\int f\,\dd\mu_Y$; $\E$ averages over the data, which physicists call the disorder.
\end{definition}

Physicists write the Gibbs weight $e^{-\beta\mathcal H}$ with an energy $\mathcal H$, so our $H_N$ is $-\beta\mathcal H$ and their free energy per variable is $f=-\phi/\beta$.\sidenote{Mathematicians call $\phi_N$ the pressure. Models from physics sum over $\{\pm1\}^N$ with the counting measure, as in \eqref{eq:freeentropy:z2post}, which adds $\log2$ to $\phi_N$.} Since $Y$ has the density $Z_N$ with respect to $Q$,
\begin{equation}\label{eq:freeentropy:mi}
\frac1NI(X;Y)=\frac1N\E H_N(X;Y)-\phi_N .
\end{equation}
So the free entropy is minus the mutual information per variable, up to the mean log-likelihood ratio of the truth, which is usually explicit. For \eqref{eq:freeentropy:z2}, with the entries above the diagonal as data and $Q$ their law under $W$, $H_N(x;A)=\lambda\sum_{i<j}A_{ij}x_ix_j-\lambda^2(N-1)/4$ has mean $\lambda^2(N-1)/4$ at the truth, and the evidence is $2^{-N}e^{-\lambda^2(N-1)/4}Z_N(A)$. \term{Replicas}\index{replicas} are independent samples $x^1,x^2,\dots$ from the same Gibbs measure, given the same data; the \term{overlap}\index{overlap} of two replicas is $R_{12}=\frac1N\ip{x^1}{x^2}$, that of a replica with the truth is $m=\frac1N\ip x{x^*}$, and a typical value $q$ of a concentrating $R_{12}$ is an \term{order parameter}\index{order parameter}. The \term{quenched}\index{quenched average} free entropy $\phi_N=\frac1N\E\log Z_N$ keeps the disorder frozen, as in a data set; the \term{annealed}\index{annealed average} free entropy $\frac1N\log\E Z_N$ averages it together with the spins.

Is $\frac1N\log Z_N$ close to its mean for typical data, what is its limit, and how do the two languages compute it? Physicists answer with confidence.

\begin{physicsbox}[self-averaging and the saddle point]
(i) The free entropy is \emph{self-averaging}: for almost every realization of the disorder, $\frac1N\log Z_N$ has the limit $\phi$. (ii) The limit is the value of a variational problem in a few order parameters at a saddle point; for the Curie--Weiss model $Z_N=\sum_\sigma\exp(\frac\beta{2N}(\sum_i\sigma_i)^2+h\sum_i\sigma_i)$ with $\beta>0$,
\[
\phi(\beta,h)=\max_{\mu\in\R}\Bigl[\log2\cosh(\beta\mu+h)-\frac{\beta\mu^2}2\Bigr],
\]
attained where $\mu=\tanh(\beta\mu+h)$. (iii) With disorder, the \emph{replica trick} gives $\phi=\lim_{n\to0}\frac1n\lim_{N\to\infty}\frac1N\log\E Z_N^n$, with $\E Z_N^n$ computed for integer $n$ by means of $n$ replicas and continued analytically to $n\to0$.
\end{physicsbox}

How do physicists derive (i)--(iii), and what is proved? We shall see that (ii) is a theorem with an explicit error at every $N$ (\cref{thm:freeentropy:curieweiss}), that (i) is one with an explicit tail at every $N$ once the mean converges (\cref{thm:freeentropy:concentration}, \cref{ch:sk}), that the integer moments in (iii) are a theorem too (\cref{prop:freeentropy:moments}), and that the step from integer $n$ to $n\to0$ is not, and cannot be in general.

\section{The physics: Hubbard--Stratonovich, saddle points and replicas}\label{sec:freeentropy:physics}

\subsection{The Curie--Weiss model by a saddle point}\label{sec:freeentropy:cw}

The \term{Curie--Weiss model}\index{Curie--Weiss model} has the Hamiltonian $H_N(\sigma)=\frac\beta{2N}S^2+hS$ on $\{\pm1\}^N$, with $S=\sum_i\sigma_i$: each spin is coupled to every other with strength $\beta/N$ and feels a field $h$, and the weight depends on $\sigma$ only through the magnetization $m=S/N$.

\begin{heuristic}[The saddle point]
The Gaussian identity $e^{a^2/2}=\int_\R e^{-u^2/2+au}\,\dd u/\sqrt{2\pi}$, with $a=\sqrt{\beta N}\,m$ and $u=\sqrt{\beta N}\,\mu$, trades the coupling of all spins for an auxiliary field $\mu$ acting on each spin separately, the \term{Hubbard--Stratonovich transformation}\index{Hubbard--Stratonovich transformation} \parencite{Stratonovich1957,Hubbard1959}. With $G(\mu)=\log2\cosh(\beta\mu+h)-\frac\beta2\mu^2$, exchanging the sum and the integral decouples the spins:
\begin{equation}\label{eq:freeentropy:hs}
Z_N=\sqrt{\frac{\beta N}{2\pi}}\int_\R e^{-N\beta\mu^2/2}\sum_\sigma e^{(\beta\mu+h)S}\dd\mu=\sqrt{\frac{\beta N}{2\pi}}\int_\R e^{NG(\mu)}\dd\mu .
\end{equation}
For large $N$ the integral is dominated by the maximum of $G$, the \term{saddle point}\index{saddle point}, where $\mu=\tanh(\beta\mu+h)$. Expanding $G$ to second order around each of its $k$ maximizers $\mu^*$, with $G''(\mu^*)=-\beta(1-\beta(1-\mu^{*2}))$, and integrating the Gaussians,
\[
\frac1N\log Z_N=G(\mu^*)+\frac1N\log\frac k{\sqrt{1-\beta(1-\mu^{*2})}}+o\Bigl(\frac1N\Bigr).
\]
The field derivative of the free entropy is the magnetization, $\gibbs m=\partial_h\phi=\tanh(\beta\mu^*+h)=\mu^*$. At $h=0$, $m=\tanh(\beta m)$ has only the solution $0$ for $\beta\le1$; for $\beta>1$ the maxima move to $\pm m^*(\beta)$, and the ferromagnet magnetizes spontaneously.
\end{heuristic}

Two steps are uncontrolled. The saddle-point evaluation needs a nondegenerate maximum; at $\beta=1$, $h=0$ the maximum of $G$ is quartic and the correction $-\frac12\log(1-\beta)$ diverges (\cref{sec:freeentropy:measuring} measures what happens instead). And reading $\gibbs m$ off the saddle point exchanges the derivative in $h$ with $N\to\infty$; at $h=0$, $\beta>1$ the exchange fails outright, since $\gibbs m=0$ at every $N$ by symmetry.

\subsection{Quenched and annealed averages}\label{sec:freeentropy:quenched}

For SK, $H_N(\sigma)=\frac\beta{\sqrt N}\sum_{i<j}g_{ij}\sigma_i\sigma_j+h\sum_i\sigma_i$ with independent standard Gaussian $g_{ij}$. At $h=0$ each $H_N(\sigma)$ is centered Gaussian with variance $\frac{\beta^2}N\binom N2$, so the annealed free entropy takes one line:
\begin{equation}\label{eq:freeentropy:annealed}
\frac1N\log\E Z_N=\log2+\frac{\beta^2(N-1)}{4N}\longrightarrow\log2+\frac{\beta^2}4 .
\end{equation}
Physicists hold that only the quenched average is physical, because $\frac1N\log Z_N$ is \term{self-averaging}\index{self-averaging}: cut a large system into blocks, and the free entropy is nearly a sum of independent block contributions, with fluctuations of order $\sqrt N$ against a mean of order $N$. (A mean-field model has no blocks; the belief rests instead on each coupling being small.) The annealed average can be dominated by rare samples: if $\log Z_N$ has Gaussian fluctuations of order $\sqrt N$, an excess $t$ over its typical value has probability about $e^{-t^2/(cN)}$ and weight $e^t$ in $\E Z_N$, and the product peaks at $t=cN/2$, of order $N$.

\subsection{The replica trick}\label{sec:freeentropy:replicas}

\begin{heuristic}[The replica trick]
For a positive random variable, $\E\log Z=\lim_{n\to0}(\E Z^n-1)/n$. For integer $n$, $Z^n$ is the partition function of $n$ replicas with the same couplings, and the Gaussian average is explicit, $\E\exp(\sum_aH_N(\sigma^a))=\exp(\frac12\Var\sum_aH_N(\sigma^a))$, with $\Cov(H_N(\sigma),H_N(\tau))=\frac{\beta^2}2(NR_{\sigma\tau}^2-1)$. For SK at $h=0$, with $R_{ab}=\frac1N\sum_i\sigma_i^a\sigma_i^b$,
\begin{equation}\label{eq:freeentropy:replicamoment}
\E Z_N^n=e^{\frac{nN\beta^2}4-\frac{n^2\beta^2}4}\sum_{\sigma^1,\dots,\sigma^n}\exp\Bigl(\frac{N\beta^2}2\sum_{a<b}R_{ab}^2\Bigr).
\end{equation}
The disorder average couples the replicas through their overlaps, and the overlap matrix becomes the order parameter. One Hubbard--Stratonovich variable $Q_{ab}$ per pair decouples the sites, leaving the $N$-th power of a sum over one column $s\in\{\pm1\}^n$: up to prefactors, $\E Z_N^n\approx\int e^{NS_n(Q)}\prod_{a<b}\dd Q_{ab}$ with
\[
S_n(Q)=\frac{n\beta^2}4-\frac{\beta^2}2\sum_{a<b}Q_{ab}^2+\log\sum_{s\in\{\pm1\}^n}e^{\beta^2\sum_{a<b}Q_{ab}s_as_b}.
\]
Take the saddle point in $Q$ as $N\to\infty$ at fixed $n$, and guess that it is \term{replica symmetric}, $Q_{ab}=q$. A last Gaussian identity gives
\[
S_n=\frac{n\beta^2}4-\frac{n(n-1)\beta^2q^2}4-\frac{n\beta^2q}2+\log\E\bigl(2\cosh(\beta\sqrt q\,Z)\bigr)^n,
\]
in which every term makes sense for real $n$. As $n\to0$, $S_n/n\to\log2+\E\log\cosh(\beta\sqrt q\,Z)+\frac{\beta^2}4(1-q)^2$, with $q=\E\tanh^2(\beta\sqrt q\,Z)$ at the stationary point: the formula of \textcite{SherringtonKirkpatrick1975}. For $\beta\le1$ the only solution is $q=0$, which gives the annealed value \eqref{eq:freeentropy:annealed}.
\end{heuristic}

Four steps are uncontrolled. The moments are computed at integers and needed near $n=0$, and integers do not determine an analytic function. The limits $N\to\infty$ and $n\to0$ are exchanged. For $n<1$ the number of pairs $n(n-1)/2$ is negative and the maximum over $Q$ becomes a minimum, a rule with no counterpart in Laplace's method. And the ansatz is a guess: its entropy at zero temperature is $-1/(2\pi)<0$, as \textcite{SherringtonKirkpatrick1975} noticed, and for $\beta>1$ the answer needs Parisi's broken replica symmetry \parencite{Parisi1979b} (\cref{ch:sk}). At $n=2$ the computation is exact.

\begin{computation}{Two replicas}\label{comp:freeentropy:secondmoment}
\step By \eqref{eq:freeentropy:replicamoment} and $(\E Z_N)^2=4^Ne^{\beta^2(N-1)/2}$, with $\rho_i=\sigma_i\tau_i$ independent uniform signs,
\[
\frac{\E Z_N^2}{(\E Z_N)^2}=e^{-\beta^2/2}\,2^{-N}\sum_\rho\exp\Bigl(\frac{\beta^2}{2N}\Bigl(\sum_i\rho_i\Bigr)^2\Bigr):
\]
the second moment of SK is a Curie--Weiss partition function at inverse temperature $\beta^2$.
\step By \cref{thm:freeentropy:curieweiss,prop:freeentropy:meanfield} below, $\frac1N\log\frac{\E Z_N^2}{(\E Z_N)^2}\to\phi(\beta^2,0)-\log2$, which is $0$ for $\beta\le1$ and positive for $\beta>1$. For $\beta<1$, Laplace's method with the prefactor (\cref{thm:toolkit:laplace}~(ii)), applied to \eqref{eq:freeentropy:hs} at inverse temperature $\beta^2$, where $G$ has the unique nondegenerate maximizer $0$, gives the limit $e^{-\beta^2/2}(1-\beta^2)^{-1/2}$.
\step If $\log Z_N-\log\E Z_N$ were Gaussian with variance $\sigma^2$, the ratio would be $e^{\sigma^2}$: so $\sigma^2(\beta)=-\frac12\log(1-\beta^2)-\frac{\beta^2}2$, the variance that \textcite{AizenmanLebowitzRuelle1987} proved for $\beta<1$, $h=0$. Numerically (\texttt{01-freeentropy-curieweiss.py}), at $\beta=0.8$ the ratio is $1.1662$ at $N=16$ and $1.2056$ at $N=200$ (limit $1.2102$); at $\beta=1.5$ it is $842$ and $2.3\times10^{38}$.
\end{computation}

The posterior \eqref{eq:freeentropy:z2post} is a special Gibbs measure, at the temperature that Bayes' rule dictates: in \eqref{eq:freeentropy:split} the ferromagnetic coupling is the square of the random one. \Textcite{Nishimori1981} singled out this line for random Ising models with a gauge argument,\sidenote{The transformation $J_{ij}\mapsto J_{ij}\tau_i\tau_j$, $\sigma_i\mapsto\sigma_i\tau_i$ leaves the Hamiltonian unchanged, and averaging over $\tau$ yields exact identities on the line, among them a closed formula for the internal energy.} and \textcite{Iba1999} read it as the line on which the Gibbs measure is the Bayesian posterior. There two replicas overlap as a replica and the truth do, $P(q)=P(m)$ in the physicist's notation, since the truth cannot be told apart from a posterior sample, and the derivative of the free entropy in the signal strength is an overlap; both statements are exact (\cref{lem:freeentropy:nishimori,thm:freeentropy:immse}).

\section{The stone}\label{sec:freeentropy:stone}

Each row pairs a step of the derivation above with the step of the proofs below that does the same work.

\begin{stone}{saddle points and replicas}{large deviations and concentration}
\stonestep{Count or linearize}{Hubbard--Stratonovich: $e^{\beta Nm^2/2}$ is a Gaussian integral over a field $\mu$; the spins decouple.}{Group configurations by magnetization, $N+1$ classes of sizes $\binom Nk\le e^{Ns(m)}$; $\mu$ returns as the Legendre dual of $m$ (\cref{lem:freeentropy:legendre}).}
\stonestep{Saddle point}{The integral is dominated by the maximum of $G$, where $\mu=\tanh(\beta\mu+h)$; Gaussian fluctuations give the $1/N$ term.}{Laplace's method on types (\cref{lem:freeentropy:types}): classes times the largest term above, a product measure below; error at most $\log(N+1)/N$, with no nondegeneracy needed.}
\stonestep{Order parameter}{Read $\gibbs m=\partial_h\phi=\mu^*$ off the saddle point; at $h=0$, $\beta>1$, pick one of two saddles.}{Derivatives of convex functions converge where the limit is differentiable, here where the maximizer is unique (\cref{cor:freeentropy:magnetization}); at $h=0$, $\beta>1$ the limit has a corner.}
\stonestep{Self-averaging}{An extensive quantity is a sum of nearly independent block contributions, with fluctuations $O(\sqrt N)$.}{$\log Z_N$ is Lipschitz in the Gaussian couplings, with constant $\beta\sqrt{(N-1)/2}$ for SK; a covariance identity gives variance and tails at every temperature (\cref{thm:freeentropy:concentration}).}
\stonestep{Quenched, not annealed}{Average $\log Z$, not $Z$: the annealed average is dominated by rare samples, right only at high temperature.}{Jensen: $\phi_N\le\frac1N\log\E Z_N$. Equality in the limit needs a second moment, for SK a Curie--Weiss partition function at $\beta^2$ (\cref{comp:freeentropy:secondmoment}).}
\stonestep{Replicate}{Compute $\E Z_N^n$ at integer $n$; the replicas couple through $R_{ab}$, and a saddle point over the overlap matrix gives the limit.}{Laplace's method on the $2^n$ column types proves the limit (\cref{prop:freeentropy:moments}); at real $n$ near $0$, $\abs{\frac1{nN}\log\E Z_N^n-\phi_N}\le\abs n\beta^2/4$ for every $N$ (\cref{prop:freeentropy:smalln}).}
\stonestep{Continue to $n\to0$}{Continue the integer-$n$ answer under an ansatz, and let $n\to0$ after $N\to\infty$.}{No counterpart: the proof never continues in $n$. Integer moments need not determine $\E\log Z$ \parencite{Tanaka2007}; for SK they grow like $e^{cNn^2}$, and Carleman's criterion fails.}
\stonestep{Bayes-optimal posteriors}{On Nishimori's line a gauge transformation gives exact identities: $P(q)=P(m)$; the derivative of the free entropy is an overlap.}{Bayes' rule: given the data, the truth is one more posterior sample (\cref{lem:freeentropy:nishimori}); with integration by parts, $I'(\gamma)=\frac12\freeentropymmse(\gamma)$ (\cref{thm:freeentropy:immse}).}
\end{stone}

Read down the right column, the stone shows what the proofs justify: the saddle point, completely and with bounds at every $N$; self-averaging, by an argument unlike the physicist's that holds at low temperature too; the integer moments; and the replica limit at real $n$ near $0$. The routes part at row~7: the physicist needs the moments near $n=0$ and has them only at integers, so the continuation is where the derivation stops being a computation and becomes a guess, refined by an ansatz and checked by its consequences. The proof does without it, at a price: to obtain the SK free entropy at low temperature (\cref{ch:sk}) it needs an idea, interpolation, that the replica computation does not contain. In the last row the routes nearly coincide: the gauge argument and Bayes' rule are one symmetry seen from two sides.

\section{The proof: Laplace's method and Gaussian concentration}\label{sec:freeentropy:laplace}

\subsection{The Curie--Weiss model}

For a probability vector $p$ on a finite set, $\freeentropyEnt(p)=-\sum_ap_a\log p_a$ (with $0\log0=0$). For a word $x\in\cA^N$ over a finite alphabet $\cA$, the \term{type}\index{method of types} $\hat p_x$ is the empirical distribution of its letters.

\begin{lemma}[Laplace's method on types]\label[lemma]{lem:freeentropy:types}
Let $\cA$ have $k\ge2$ elements, $\Delta$ be the set of probability vectors on $\cA$, and $F\colon\Delta\to\R$ be convex and continuous. With $\Lambda=\max_{p\in\Delta}[\freeentropyEnt(p)+F(p)]$, for every $N\ge1$,
\[
\Lambda\le\frac1N\log\sum_{x\in\cA^N}e^{NF(\hat p_x)}\le\Lambda+\frac{(k-1)\log(N+1)}N.
\]
\end{lemma}

\begin{proof}
The plan is to group the words by type for the upper bound and to use the Gibbs variational principle with an i.i.d.\ measure for the lower bound. A type is fixed by the counts of $k-1$ letters, so there are at most $(N+1)^{k-1}$ types. Each word of type $p$ has probability $\prod_ap_a^{Np_a}=e^{-N\freeentropyEnt(p)}$ under $p^{\otimes N}$, so at most $e^{N\freeentropyEnt(p)}$ words have type $p$, and the sum is at most $(N+1)^{k-1}e^{N\Lambda}$. For the lower bound, Jensen's inequality for the concave logarithm gives, for every probability vector $\nu$ on a finite set and every function $H$ on it, $\log\sum e^H\ge\log\sum_{\nu>0}\nu\,e^H/\nu\ge\sum\nu H+\freeentropyEnt(\nu)$ (the Gibbs variational principle). Take $\nu=p^{\otimes N}$, of entropy $N\freeentropyEnt(p)$, and $H(x)=NF(\hat p_x)$. Under $\nu$ the type has mean $p$, so by convexity $\sum_x\nu(x)F(\hat p_x)\ge F(p)$ and $\log\sum_xe^{NF(\hat p_x)}\ge N(\freeentropyEnt(p)+F(p))$. Maximize over $p$.
\end{proof}

Write $s(m)=-\frac{1+m}2\log\frac{1+m}2-\frac{1-m}2\log\frac{1-m}2$ for the entropy of a $\pm1$ variable with mean $m\in[-1,1]$.\sidenote{The limit, without the bounds, also follows from Varadhan's lemma (\cref{thm:toolkit:laplace}~(iii)): by Cramér's theorem the type of $N$ independent uniform letters satisfies the large deviation principle with the good rate function $\log k-\freeentropyEnt(p)$, and $F$ is continuous and bounded above on the simplex. Convexity of $F$ buys the exact lower bound.}

\begin{lemma}[Legendre duality]\label[lemma]{lem:freeentropy:legendre}
For every $y\in\R$, $\max_{m\in[-1,1]}[s(m)+ym]=\log2\cosh y$, attained only at $m=\tanh y$.
\end{lemma}

\begin{proof}
On $(-1,1)$, $s'(m)=-\operatorname{artanh}m$, which tends to $\mp\infty$ at $\pm1$, and $s$ is strictly concave, so $s(m)+ym$ has the unique maximizer $\tanh y$. With $p=\frac{1+\tanh y}2$, $\log p=y-\log2\cosh y$ and $\log(1-p)=-y-\log2\cosh y$, so $s(\tanh y)=\log2\cosh y-y\tanh y$.
\end{proof}

\begin{theorem}[The Curie--Weiss free entropy]\label{thm:freeentropy:curieweiss}
Let $\beta\ge0$, $h\in\R$, $Z_N(\beta,h)=\sum_{\sigma\in\{\pm1\}^N}e^{\frac\beta{2N}S^2+hS}$ with $S=\sum_i\sigma_i$, and
\[
\phi(\beta,h)=\max_{m\in[-1,1]}\Phi(m),\qquad\Phi(m)=s(m)+\frac\beta2m^2+hm.
\]
\begin{enumerate}
\item For every $N\ge1$, $\phi(\beta,h)\le\frac1N\log Z_N(\beta,h)\le\phi(\beta,h)+\frac{\log(N+1)}N$.
\item If $\beta>0$, then $\phi(\beta,h)=\max_{\mu\in\R}G(\mu)$ with $G(\mu)=\log2\cosh(\beta\mu+h)-\frac\beta2\mu^2$; $\Phi$ and $G$ have the same maximizers, each in $(-1,1)$ and a solution of $m=\tanh(\beta m+h)$.
\end{enumerate}
\end{theorem}

\begin{proof}
The plan is to read part~1 off \cref{lem:freeentropy:types} and part~2 off a function of both $\mu$ and $m$. With $\cA=\{\pm1\}$ the type of $\sigma$ is $(\frac{1+m}2,\frac{1-m}2)$, $m=S/N$, of entropy $s(m)$, and $\frac\beta2m^2+hm$ is convex in the type because $m$ is linear in it and $\beta\ge0$: this is part~1. For part~2 let $J(\mu,m)=s(m)+(\beta\mu+h)m-\frac\beta2\mu^2$. For fixed $m$ it is maximized only at $\mu=m$, with value $\Phi(m)$; for fixed $\mu$ only at $m=\tanh(\beta\mu+h)$, with value $G(\mu)$ (\cref{lem:freeentropy:legendre}). Both maxima exist ($\Phi$ is continuous on a compact set, $G\to-\infty$ at $\pm\infty$) and equal $\sup J$. If $m^*$ maximizes $\Phi$, then $J(m^*,m^*)=\sup J$, so $m^*$ maximizes $G$, and uniqueness in $m$ gives $m^*=\tanh(\beta m^*+h)$. If $\mu^*$ maximizes $G$ and $m=\tanh(\beta\mu^*+h)$, then $J(\mu^*,m)=\sup J$, so $m$ maximizes $\Phi$, and uniqueness in $\mu$ gives $\mu^*=m$.
\end{proof}

So physics box (ii) holds, with error at most $\log(N+1)/N$ at every $N$, whatever the nature of the maximum; the lower bound is the free entropy of the best product measure, the \term{mean-field bound}\index{mean-field bound}.

\begin{proposition}[The mean-field equation]\label[proposition]{prop:freeentropy:meanfield}
Let $\beta\ge0$. If $h\ne0$, $\Phi$ has a unique maximizer $m^*(\beta,h)$, of the sign of $h$. If $h=0$ and $\beta\le1$, the unique maximizer is $0$ and $\phi=\log2$. If $h=0$ and $\beta>1$, the maximizers are $\pm m^*(\beta)$, with $m^*(\beta)$ the unique solution in $(0,1)$ of $m=\tanh(\beta m)$.
\end{proposition}

\begin{proof}
The plan is to count the solutions of $m=\tanh(\beta m+h)$, which contain the maximizers by \cref{thm:freeentropy:curieweiss} (for $\beta=0$, $\Phi$ is strictly concave with the single maximizer $\tanh h$). A concave $u$ on $[0,1]$ with $u(0)>0$, or strictly concave with $u(0)=0$, has at most one zero in $(0,1]$: zeros $0<z_1<z_2$ would give $u(z_1)\ge(1-\frac{z_1}{z_2})u(0)+\frac{z_1}{z_2}u(z_2)$, strictly in the second case, so $u(z_1)>0$. If $h>0$ (the case $h<0$ is symmetric), $\Phi(m)>\Phi(-m)$ for $m>0$ and $\Phi'(0)=h>0$, so the maximizers are zeros in $(0,1)$ of $u(m)=\tanh(\beta m+h)-m$, concave on $[0,1]$ with $u(0)>0>u(1)$: there is exactly one. If $h=0$, $\beta\le1$ and $m\ne0$, then $\abs{\tanh(\beta m)}<\abs m$, so only $0$ remains. If $h=0$, $\beta>1$, then $u(m)=\tanh(\beta m)-m$ is strictly concave with $u(0)=0$, $u'(0)>0>u(1)$, so it has one zero $m^*$ in $(0,1)$; the solution $0$ is a local minimum since $\Phi''(0)=\beta-1>0$, and $\Phi$ is even.
\end{proof}

\begin{corollary}[Magnetization]\label[corollary]{cor:freeentropy:magnetization}
Let $\beta\ge0$ and $h\ne0$. Then the Gibbs average $\gibbs m_N=\gibbs{S/N}$ converges to $m^*(\beta,h)$.
\end{corollary}

\begin{proof}
The plan is to differentiate in $h$ through convexity. The function $f_N(h)=\frac1N\log Z_N(\beta,h)$ is convex with $f_N'(h)=\gibbs m_N$, and $f_N\to\phi$. Since $\frac{f_N(h)-f_N(h-\delta)}\delta\le f_N'(h)\le\frac{f_N(h+\delta)-f_N(h)}\delta$ for $\delta>0$, it suffices that $\phi$ is differentiable at $h$ with derivative $m^*$. Write $\phi(h)=\max_m[\Phi_0(m)+hm]$ with $\Phi_0=s+\frac\beta2m^2$. For $h'\ne h$ with a maximizer $m'$, comparing values gives $\phi(h')-\phi(h)\ge(h'-h)m^*$ and $\phi(h)-\phi(h')\ge(h-h')m'$, so the difference quotient lies between $m^*$ and $m'$. As $h'\to h$, every limit point of $m'$ maximizes $\Phi_0(m)+hm$, because $\phi$ (convex) and $\Phi_0$ are continuous, so it equals $m^*$ (\cref{prop:freeentropy:meanfield}).
\end{proof}

At $h=0$, $\beta>1$ the same comparison gives $\partial_h\phi(\beta,0\pm)=\pm m^*(\beta)$: the limit has a corner while $\gibbs m_N=0$, and choosing a saddle point is choosing the order of the limits $h\to0^\pm$ and $N\to\infty$. The lemma on types also proves the integer-moment step of physics box (iii).

\begin{proposition}[Integer moments of SK]\label[proposition]{prop:freeentropy:moments}
Let $h=0$, $\beta\ge0$, $n\ge1$ an integer and $N\ge2$. Let $\Lambda_n(\beta)$ be the maximum of $\freeentropyEnt(p)+\frac{\beta^2}2\sum_{a<b}(\E_ps_as_b)^2$ over probability vectors $p$ on $\{\pm1\}^n$, with $s\sim p$. Then
\[
0\le\frac1N\log\E Z_N^n-\frac{n\beta^2}4+\frac{n^2\beta^2}{4N}-\Lambda_n(\beta)\le\frac{(2^n-1)\log(N+1)}N.
\]
\end{proposition}

\begin{proof}
Formula \eqref{eq:freeentropy:replicamoment} is exact, since it uses only the Gaussian moment generating function. Let $\cA=\{\pm1\}^n$, and let $x_i=(\sigma^1_i,\dots,\sigma^n_i)$ be the columns of the replicas. Then $R_{ab}=\E_{\hat p_x}s_as_b$ is linear in the type, so the function $F(p)$ in the definition of $\Lambda_n$, without the entropy, is convex, and the sum in \eqref{eq:freeentropy:replicamoment} is $\sum_xe^{NF(\hat p_x)}$; apply \cref{lem:freeentropy:types} with $k=2^n$.
\end{proof}

For $n=1$, $\Lambda_1=\log2$ gives the annealed value \eqref{eq:freeentropy:annealed}; for $n=2$ the maximum reduces to the law of $s_1s_2$, and $\Lambda_2=\log2+\phi(\beta^2,0)$, which is \cref{comp:freeentropy:secondmoment}. Exchanging maxima as in the proof of \cref{thm:freeentropy:curieweiss}, with $\frac{\beta^2}2r^2=\max_Q[\beta^2Qr-\frac{\beta^2}2Q^2]$ for each pair, gives $\frac{n\beta^2}4+\Lambda_n=\max_QS_n(Q)$: at integer $n$ the physicist's saddle point is a maximum, and it is exact.

\subsection{Gaussian concentration}\label{sec:freeentropy:concentration}

A \term{Gaussian disordered system} is a Gibbs measure on a finite set $\Sigma$ with weights $\pi(\sigma)>0$ and a Hamiltonian $H(\sigma)=\ip g{a(\sigma)}+b(\sigma)$ linear in a standard Gaussian vector $g\in\R^K$, so that $Z(g)=\sum_\sigma\pi(\sigma)e^{H(\sigma)}$; $a$ and $b$ are deterministic, or independent of $g$ and then everything holds conditionally on them. Let $L=\max_\sigma\norm{a(\sigma)}$. For SK, $a_{ij}(\sigma)=\frac\beta{\sqrt N}\sigma_i\sigma_j$ ($i<j$) and $L^2=\frac{\beta^2(N-1)}2$; the posterior \eqref{eq:freeentropy:z2post} has the same $L$ with $\beta=\lambda$, given $x^*$.

The tool is the covariance identity of the appendix (\cref{lem:toolkit:covariance}): for independent standard Gaussian vectors $X,Y$, with $X_t=tX+\sqrt{1-t^2}\,Y$, and twice continuously differentiable $f_1,f_2$ which, with their first and second derivatives, have moderate growth,
\begin{equation}\label{eq:freeentropy:covariance}
\Cov\bigl(f_1(X),f_2(X)\bigr)=\int_0^1\E\ip{\nabla f_1(X)}{\nabla f_2(X_t)}\dd t.
\end{equation}

\begin{theorem}[Concentration of the free entropy]\label{thm:freeentropy:concentration}
For a Gaussian disordered system, $F=\log Z$ satisfies, for every $\theta\in\R$ and $t\ge0$,
\begin{enumerate}
\item $\Var F\le\E\norm{\gibbs a}^2\le L^2$;
\item $\E e^{\theta(F-\E F)}\le e^{\theta^2L^2/2}$;
\item $\Prob(\abs{F-\E F}\ge t)\le2e^{-t^2/(2L^2)}$ if $L>0$.
\end{enumerate}
For SK, with every $\beta>0$, $h\in\R$, $N\ge2$ and $u\ge0$,
\begin{gather}
\Var\log Z_N\le\frac{\beta^2}2\bigl(N\E\gibbs{R_{12}^2}-1\bigr)\le\frac{\beta^2(N-1)}2,\notag\\
\Prob\Bigl(\Bigl|\frac1N\log Z_N-\phi_N\Bigr|\ge u\Bigr)\le2e^{-\frac{N^2u^2}{\beta^2(N-1)}}\le2e^{-Nu^2/\beta^2}.\label{eq:freeentropy:sktail}
\end{gather}
\end{theorem}

\begin{proof}
The plan is to apply \eqref{eq:freeentropy:covariance} with $f_1=f_2=F$ for the variance and with $f_1=F$, $f_2=e^{\theta F}$ for a differential inequality on the moment generating function. Since $\partial_kF=\gibbs{a_k}$, $\norm{\nabla F}=\norm{\gibbs a}\le\gibbs{\norm a}\le L$, and the Hessian of $F$, the Gibbs covariance matrix of $a$, is bounded; so $\abs{F(x)}\le\abs{F(0)}+L\norm x$, and $F$ and $e^{\theta F}$ satisfy the growth conditions.

1. This is the Gaussian Poincaré inequality (\cref{thm:toolkit:poincare}), which \eqref{eq:freeentropy:covariance} with $f_1=f_2=F$ and Cauchy--Schwarz give at once, since $X_t\sim X$: $\Var F\le\E\norm{\nabla F}^2$.

2. Let $\bar F=F-\E F$ and $M(\theta)=\E e^{\theta\bar F}$. Since $\E\bar F=0$, $M'(\theta)$ equals $\Cov(F,e^{\theta\bar F})$. The gradient of $e^{\theta\bar F}$ is $\theta e^{\theta\bar F}\nabla F$, and $\ip{\nabla F(X)}{\nabla F(X_t)}\le L^2$, so for $\theta\ge0$
\[
M'(\theta)=\theta\int_0^1\E\bigl[e^{\theta\bar F(X_t)}\ip{\nabla F(X)}{\nabla F(X_t)}\bigr]\dd t\le\theta L^2M(\theta).
\]
So $(\log M)'\le\theta L^2$ on $[0,\infty)$, and $M(0)=1$ gives $\log M(\theta)\le\theta^2L^2/2$. The argument used only $\norm{\nabla F}\le L$ and the growth conditions, which $-F$ satisfies too; this covers $\theta<0$.

3. Markov's inequality gives $\Prob(\bar F\ge t)\le e^{-\theta t+\theta^2L^2/2}$, which is $e^{-t^2/(2L^2)}$ at $\theta=t/L^2$; the same holds for $-\bar F$.

For SK, $\norm{\gibbs a}^2=\frac{\beta^2}N\sum_{i<j}\gibbs{\sigma_i\sigma_j}^2$, and the square of a Gibbs average is an average over two replicas: $\sum_{i<j}\gibbs{\sigma_i\sigma_j}^2=\gibbs{\sum_{i<j}\sigma^1_i\sigma^1_j\sigma^2_i\sigma^2_j}=\frac12\gibbs{N^2R_{12}^2-N}$. Since $R_{12}^2\le1$, this gives both variance bounds; the tail is part~3 with $t=Nu$.
\end{proof}

The tail \eqref{eq:freeentropy:sktail} is summable in $N$, so by the Borel--Cantelli lemma $\frac1N\log Z_N-\phi_N\to0$ almost surely, whatever the coupling of the disorders of different sizes. Since $\phi_N$ converges for SK at every $\beta$ and $h$ (\cref{thm:sk:guerratoninelli}), physics box (i) is a theorem.\sidenote{Part~1 is a statement about replicas: the variance of $\log Z_N$ is controlled by the mean square overlap, which is of order $1/N$ when the overlap concentrates at $0$, as at high temperature.} Part~2 settles the replica limit at real $n$.

\begin{proposition}[The replica free entropy at small $n$]\label[proposition]{prop:freeentropy:smalln}
For a Gaussian disordered system and every real $n\ne0$, $\frac1n\log\E Z^n-\E\log Z$ lies between $0$ and $\frac{nL^2}2$. In particular $\E\log Z=\lim_{n\to0}\frac1n\log\E Z^n=\lim_{n\to0}\frac{\E Z^n-1}n$, and for SK $\abs{\frac1{nN}\log\E Z_N^n-\phi_N}\le\frac{\abs n\beta^2}4$ for every $N$.
\end{proposition}

\begin{proof}
In the notation of the last proof, $\frac1n\log\E Z^n-\E\log Z=\frac1n\log M(n)$, and $0\le\log M(n)\le\frac{n^2L^2}2$ by Jensen's inequality ($M(n)\ge e^{n\E\bar F}=1$) and part~2. Since $\log\E Z^n\to0$, $\frac{\E Z^n-1}n=\frac{e^{\log\E Z^n}-1}{\log\E Z^n}\cdot\frac{\log\E Z^n}n\to\E\log Z$. For SK, $L^2/N\le\beta^2/2$.
\end{proof}

The proposition is the rigorous content of the replica limit: at real $n$ near $0$ the replica free entropy is the quenched one, uniformly in $N$, so the limits $n\to0$ and $N\to\infty$ commute. It gives no way to obtain $\E Z_N^n$ at non-integer $n$ from its values at integers. For integer $n\ge1$, $Z_N^n\ge\sum_\sigma e^{nH_N(\sigma)}$, so for SK at $h=0$, $\E Z_N^n\ge2^Ne^{n^2\beta^2(N-1)/4}$. The moments grow like $e^{cn^2}$, so $\sum_n(\E Z_N^n)^{-1/(2n)}<\infty$, and \term{Carleman's condition}\index{Carleman's condition}, the standard sufficient condition for moments to determine a law on $[0,\infty)$ \parencite{Schmudgen2017}, fails at every $N\ge2$, $\beta>0$. Such growth is the signature of the lognormal law, which its moments do not determine \parencite{Heyde1963} (\cref{ex:freeentropy:heyde}). \Textcite{Tanaka2007} built counterexamples of this kind, among them a random two-spin Ising model, in which the moments of $Z$ do not determine $\E\log Z$, and concluded that the replica method cannot be justified unconditionally.

\section{The proof: the scalar channel, the Nishimori identity and I-MMSE}\label{sec:freeentropy:channel}

The simplest posterior with a planted truth is that of one variable seen through Gaussian noise; the replica-symmetric formulas of Part~II reduce high-dimensional posteriors to it, and the two identities proved here drive the proofs there.

Let $X\sim P_0$ with $\E X^2<\infty$, $Z\sim\Normal(0,1)$ independent, and $Y=\sqrt\gamma\,X+Z$ with $\gamma\ge0$. Against $Q=\Normal(0,1)$ the log-likelihood ratio is $\sqrt\gamma\,xy-\frac\gamma2x^2$, so in the notation of \cref{def:freeentropy:gibbs}
\[
\mu_y(\dd x)=\frac{e^{\sqrt\gamma xy-\gamma x^2/2}}{\cZ_\gamma(y)}P_0(\dd x),\qquad\psi(\gamma)=\E\log\cZ_\gamma(Y),
\]
with $\cZ_\gamma(y)=\int e^{\sqrt\gamma xy-\gamma x^2/2}P_0(\dd x)$. The posterior mean $\gibbs x$ minimizes the mean squared error, $\freeentropymmse(\gamma)=\E(X-\gibbs x)^2$, and with $I(\gamma)=I(X;Y)$, \eqref{eq:freeentropy:mi} reads
\begin{equation}\label{eq:freeentropy:scalarmi}
I(\gamma)=\frac\gamma2\E X^2-\psi(\gamma).
\end{equation}
For $N$ independent copies the posterior is a product and the free entropy per variable is $\psi(\gamma)$. For the Rademacher prior $\gibbs x=\tanh(\sqrt\gamma Y)$ and $\psi(\gamma)=\E\log\cosh(\gamma+\sqrt\gamma Z)-\frac\gamma2$; for the standard Gaussian prior $\gibbs x=\sqrt\gamma Y/(1+\gamma)$, $\freeentropymmse=1/(1+\gamma)$ and $I=\frac12\log(1+\gamma)$.

\begin{lemma}[Nishimori identity]\label[lemma]{lem:freeentropy:nishimori}
Let $(X,Y)$ take values in Polish spaces, let $\mu_Y$ be a regular conditional law of $X$ given $Y$, and given $(X,Y)$ let $x^1,\dots,x^k$ be independent with law $\mu_Y$. For every measurable $f\ge0$, and every $f$ for which one side is finite with $\abs f$ in place of $f$,
\[
\E\gibbs{f(Y,x^1,\dots,x^{k-1},x^k)}=\E\gibbs{f(Y,x^1,\dots,x^{k-1},X)}.
\]
\end{lemma}

\begin{proof}
The joint law of $(Y,X,x^1,\dots,x^k)$ is
\[
P_Y(\dd y)\,\mu_y(\dd x)\,\mu_y(\dd x^1)\cdots\mu_y(\dd x^k):
\]
given $Y$, the truth and the replicas are independent with the same law, so exchanging $X$ and $x^k$ leaves the joint law unchanged.
\end{proof}

\begin{corollary}[First consequences]\label[corollary]{cor:freeentropy:nishimori}
In the scalar channel, $\E\gibbs{x^2}=\E X^2$, $\E[X\gibbs x]=\E\gibbs x^2$, and $\freeentropymmse(\gamma)=\E X^2-\E\gibbs x^2$. For every prior on $\R^N$ and every data model, the overlap $R_{12}$ of two replicas under $\E\gibbs\cdot$ has the law of the overlap $m$ of a replica with the truth.
\end{corollary}

\begin{proof}
Apply \cref{lem:freeentropy:nishimori} with $k=1$, $f=x^2$, then with $k=2$, $f=x^1x^2$, which is integrable since $\E\gibbs{\abs{x^1x^2}}=\E\gibbs{\abs x}^2\le\E\gibbs{x^2}=\E X^2$, and expand $\E(X-\gibbs x)^2$. For the overlaps take $f=\ind\{\frac1N\ip{x^1}{x^2}\in B\}$ for Borel sets $B$.
\end{proof}

\begin{theorem}[I-MMSE]\label{thm:freeentropy:immse}
If $P_0$ has bounded support, then $I(\gamma)=\frac12\int_0^\gamma\freeentropymmse(r)\,\dd r$ for every $\gamma\ge0$. In particular $I'(\gamma)=\frac12\freeentropymmse(\gamma)$ for $\gamma>0$, and from the right at $0$; equivalently, $\psi'(\gamma)=\frac12\E\gibbs x^2$.
\end{theorem}

\begin{proof}
The plan is to substitute $Y=\sqrt\gamma X+Z$, so that the law of $(X,Z)$ does not depend on $\gamma$, to differentiate, and to remove the noise by integration by parts and the Nishimori identity. Let $\abs x\le K$ on the support of $P_0$, and $\psi(\gamma)=\E\ell(\gamma)$ with
\[
\ell(\gamma)=\log\int\exp\Bigl(\gamma Xx+\sqrt\gamma Zx-\frac\gamma2x^2\Bigr)P_0(\dd x).
\]
For fixed $(X,Z)$ and $\gamma>0$, $\ell'(\gamma)=\gibbs{Xx+\frac{Zx}{2\sqrt\gamma}-\frac{x^2}2}$, which for $\gamma\ge\gamma_0>0$ is at most $\frac32K^2+\frac{K\abs Z}{2\sqrt{\gamma_0}}$ in absolute value. By dominated convergence,
\[
\psi'(\gamma)=\E[X\gibbs x]-\frac12\E\gibbs{x^2}+\frac1{2\sqrt\gamma}\E[Z\gibbs x].
\]
Given $X$, $\gibbs x$ is a smooth function of $Z$ with $\partial_Z\gibbs x=\sqrt\gamma(\gibbs{x^2}-\gibbs x^2)$, bounded by $\sqrt\gamma K^2$, so Stein's identity (\cref{thm:toolkit:gaussianibp}) gives $\E[Z\gibbs x]=\sqrt\gamma\,\E[\gibbs{x^2}-\gibbs x^2]$. Hence $\psi'(\gamma)=\E[X\gibbs x]-\frac12\E\gibbs x^2=\frac12\E\gibbs x^2$ by \cref{cor:freeentropy:nishimori}, and by \eqref{eq:freeentropy:scalarmi} $I'(\gamma)=\frac12(\E X^2-\E\gibbs x^2)=\frac12\freeentropymmse(\gamma)$. For fixed $(X,Z)$, $\gibbs x$ and $\ell$ are continuous in $\gamma\ge0$ and dominated as above, so $\freeentropymmse$ is continuous on $[0,\infty)$ and $I$ is continuous at $0$, with $I(0)=0$. Integrate from $\epsilon$ to $\gamma$ and let $\epsilon\to0$.
\end{proof}

The Gaussian prior, which is unbounded, satisfies the theorem by the formulas above.\sidenote{For the Rademacher prior, \cref{cor:freeentropy:nishimori} gives $\freeentropymmse(\gamma)=1-\E\tanh(\gamma+\sqrt\gamma Z)$ and the identity $\E\tanh(\gamma+\sqrt\gamma Z)=\E\tanh^2(\gamma+\sqrt\gamma Z)$, which is not apparent from the formulas.} \Textcite{GuoShamaiVerdu2005} proved the relation for every input of finite variance, for vector inputs and in continuous time; their generality adds control of the posterior moments for unbounded support. The vector form $\frac{\dd}{\dd\gamma}I(X;\sqrt\gamma X+Z)=\frac12\E\norm{X-\gibbs x}^2$ has the same proof (\cref{ex:freeentropy:vector}); the interpolation proofs of \cref{ch:nishimori,ch:glm} differentiate it, and there too the derivative of a free entropy in a signal-to-noise ratio is an overlap.

\section{Measuring it}\label{sec:freeentropy:measuring}

\begin{figure}[t]
\begin{fullwidth}
\centering
\begin{tikzpicture}
\begin{groupplot}[group style={group size=3 by 1,horizontal sep=13mm},rplot,width=57mm,height=52mm,
  title style={font=\small}]
\nextgroupplot[xmode=log,ymode=log,xmin=3.6,xmax=21,ymin=0.005,ymax=0.7,
  xtick={4,8,16},xticklabels={4,8,16},xlabel={$N$},ylabel={std.\ dev.\ of $\frac1N\log Z_N$},
  title={(a) SK: self-averaging}]
\addplot[proofpred] table[x=N,y=bd15] {figures/data/01-freeentropy-sk-sd.dat};
\addplot[proofpred] table[x=N,y=bd08] {figures/data/01-freeentropy-sk-sd.dat};
\addplot[proofpred] table[x=N,y=bd05] {figures/data/01-freeentropy-sk-sd.dat};
\addplot[physpred] table[x=N,y=alr08] {figures/data/01-freeentropy-sk-sd.dat};
\addplot[physpred] table[x=N,y=alr05] {figures/data/01-freeentropy-sk-sd.dat};
\addplot[measured3,ebars] table[x=N,y=sd15,y error plus=hi15,y error minus=lo15] {figures/data/01-freeentropy-sk-sd.dat};
\addplot[measured2,ebars] table[x=N,y=sd08,y error plus=hi08,y error minus=lo08] {figures/data/01-freeentropy-sk-sd.dat};
\addplot[measured,ebars] table[x=N,y=sd05,y error plus=hi05,y error minus=lo05] {figures/data/01-freeentropy-sk-sd.dat};
\node[font=\scriptsize,text=proof,anchor=west] at (axis cs:16.4,0.256) {$1.5$};
\node[font=\scriptsize,text=proof,anchor=west] at (axis cs:16.4,0.137) {$0.8$};
\node[font=\scriptsize,text=proof,anchor=west] at (axis cs:16.4,0.0856) {$0.5$};
\nextgroupplot[xmin=0,xmax=3.1,ymin=0.7,ymax=1.9,xlabel={$n$},ylabel={$\frac1{nN}\log\E Z_N^n$},
  title={(b) SK: replica moments, $N=12$}]
\addplot[proofpred] table[x=n,y=up15] {figures/data/01-freeentropy-sk-moments.dat};
\addplot[proofpred] table[x=n,y=q15] {figures/data/01-freeentropy-sk-moments.dat};
\addplot[proofpred] table[x=n,y=up05] {figures/data/01-freeentropy-sk-moments.dat};
\addplot[proofpred] table[x=n,y=q05] {figures/data/01-freeentropy-sk-moments.dat};
\addplot[only marks,mark=diamond*,mark size=2.3pt,color=proof] table[x=n,y=ex15] {figures/data/01-freeentropy-sk-exact.dat};
\addplot[only marks,mark=diamond*,mark size=2.3pt,color=proof] table[x=n,y=ex05] {figures/data/01-freeentropy-sk-exact.dat};
\addplot[measured3,ebars] table[x=n,y=mc15,y error plus=hi15,y error minus=lo15] {figures/data/01-freeentropy-sk-moments.dat};
\addplot[measured,ebars] table[x=n,y=mc05,y error plus=hi05,y error minus=lo05] {figures/data/01-freeentropy-sk-moments.dat};
\node[font=\scriptsize,text=ink,anchor=west] at (axis cs:1.9,1.20) {$\beta=1.5$};
\node[font=\scriptsize,text=ink,anchor=west] at (axis cs:1.9,0.80) {$\beta=0.5$};
\nextgroupplot[xmode=log,xmin=8,xmax=1.3e6,ymin=-0.2,ymax=4.6,xlabel={$N$},
  xtick={1e1,1e2,1e3,1e4,1e5,1e6},ylabel={$N\bigl(\frac1N\log Z_N-\phi\bigr)$},title={(c) Curie--Weiss: saddle point}]
\addplot[proofpred] table[x=N,y=bound] {figures/data/01-freeentropy-cw.dat};
\addplot[guide] coordinates {(8,0) (1.3e6,0)};
\addplot[physpred] table[x=N,y=pred10] {figures/data/01-freeentropy-cw.dat};
\addplot[physpred] table[x=N,y=pred15] {figures/data/01-freeentropy-cw.dat};
\addplot[physpred] table[x=N,y=pred05] {figures/data/01-freeentropy-cw.dat};
\addplot[measured2] table[x=N,y=gap10] {figures/data/01-freeentropy-cw.dat};
\addplot[measured3] table[x=N,y=gap15] {figures/data/01-freeentropy-cw.dat};
\addplot[measured] table[x=N,y=gap05] {figures/data/01-freeentropy-cw.dat};
\node[font=\scriptsize,text=ink,anchor=east] at (axis cs:3e3,2.75) {$\beta=1$};
\node[font=\scriptsize,text=ink,anchor=west] at (axis cs:3e3,1.30) {$\beta=1.5$};
\node[font=\scriptsize,text=ink,anchor=west] at (axis cs:3e3,0.65) {$\beta=0.5$};
\end{groupplot}
\end{tikzpicture}
\end{fullwidth}
\caption{Exact computations against the predictions. (a) SK at $h=0$ by exact enumeration of the $2^N$ configurations: standard deviation of $\frac1N\log Z_N$ over $M$ disorder samples ($M=4000$ for $N\le8$, $3000$ for $N=10,12$, $800$ for $N=14$, $300$ for $N=16$; seed $100+N$) at $\beta=0.5$ (dots), $0.8$ (circles) and $1.5$ (squares), with $95\%$ bootstrap intervals ($1000$ resamples). Lapis dashed: the bound $\beta\sqrt{(N-1)/2}/N$ of \cref{thm:freeentropy:concentration}, labeled by $\beta$; carnelian: the high-temperature prediction $\sigma(\beta)/N$ of \cref{comp:freeentropy:secondmoment} at $\beta=0.5,0.8$. (b) Replica free entropy of SK at $N=12$ from the $3000$ samples of (a), Monte Carlo with $95\%$ bootstrap intervals ($400$ resamples); lapis diamonds: exact values at $n=1,2,3$ from \eqref{eq:freeentropy:replicamoment}; lapis dashed: the band $[\phi_N,\phi_N+n\beta^2(N-1)/(4N)]$ of \cref{prop:freeentropy:smalln}, with $\phi_N$ from the samples. (c) The Curie--Weiss model at $h=0$, exact sums over the $N+1$ magnetizations at $41$ sizes (no intervals). Lapis dashed: the bound $\log(N+1)$ of \cref{thm:freeentropy:curieweiss}, which leaves the panel at $N\approx100$; carnelian: the saddle-point predictions $-\frac12\log(1-\beta)$ at $\beta=0.5$, $\log2-\frac12\log(1-\beta(1-m^{*2}))$ (two saddles) at $\beta=1.5$, and $\frac14\log N+\log c$, $c=2\cdot12^{1/4}\Gamma(\frac54)/\sqrt{2\pi}$, at $\beta=1$. Scripts: \texttt{figures/scripts/01-freeentropy-sk.py}, \texttt{01-freeentropy-curieweiss.py}.}\label{fig:freeentropy:sk}
\end{figure}

The script \texttt{01-freeentropy-sk.py} enumerates SK at $N=4,\dots,16$ for $\beta=0.5$, $0.8$, $1.5$ at $h=0$, in about a minute. In \cref{fig:freeentropy:sk}(a), at $\beta=0.5$ the standard deviation of $\frac1N\log Z_N$ falls like $1/N$, from $0.027$ at $N=4$ to $0.0082$ at $N=16$, as the variance of $\log Z_N$ climbs toward $\sigma^2(0.5)=0.0188$ ($0.0172$ at $N=16$); at $\beta=1.5$ the variance keeps growing, from $0.53$ to $1.53$, as expected at low temperature.\sidenote{At $\beta=0.8$ the approach is slower, $0.125$ at $N=16$ against the limit $0.191$, an effect of the nearby critical point.} The bound of \cref{thm:freeentropy:concentration} holds and is crude: at $N=16$, $L^2=1.88$, $4.80$ and $16.9$ against variances of $0.017$, $0.125$ and $1.53$. The intermediate bound $\frac{\beta^2}2(N\E\gibbs{R_{12}^2}-1)$ of part~1, measured on the same samples, is $0.036$, $0.30$ and $4.5$: within a factor of three, levelling off at $\beta=0.5$ and growing roughly linearly in $N$ at $\beta=1.5$, where the overlap does not concentrate at $0$.

Panel~(b) puts the quenched and annealed averages on one curve, at $n\to0$ and $n=1$. At $\beta=0.5$ the curve is flat, from $0.7497$ to $0.7519$, and the Monte Carlo estimates agree with the exact integer moments. At $\beta=1.5$ the estimate at $n=1$, $1.213$ with interval $[1.202,1.225]$, agrees with the exact $1.209$, but at $n=2$ and $3$ the estimates $1.330$ and $1.406$ fall far below the exact $1.416$ and $1.817$, and even below the rigorous lower bounds $\frac{\log2}n+\frac{n\beta^2(N-1)}{4N}=1.378$ and $1.778$. This is the physicist's warning about rare samples, measured: $3000$ samples miss those that dominate $\E Z_N^2$, and resampling cannot reveal what was never sampled.

At $N=16$, with $95\%$ half-widths, $\phi_N=0.7517\pm0.0009$ against the annealed $0.7517$ at $\beta=0.5$, where the Gaussian fluctuations of \cref{comp:freeentropy:secondmoment}, with mean $-\sigma^2/2$ so that $Z_N/\E Z_N$ keeps mean one, predict a deficit $\sigma^2(0.5)/(2N)=0.0006$, below the resolution of this run; and $1.164\pm0.009$ against $1.220$ at $\beta=1.5$, where the limit lies below Guerra's replica-symmetric bound $1.2388$ (\cref{ch:sk}).\sidenote{The bound is attained at $q=0.3526$. The one-step replica-symmetry-breaking value of the Parisi functional, also an upper bound, is lower by $2.2\times10^{-4}$, and the annealed limit is $1.2556$.} Since $N\phi_N$ is superadditive (\cref{thm:sk:guerratoninelli}), each $\phi_N$ is a lower bound on the limit; the gap between $\phi_{16}$ and the one-step value, $0.075$, is over three hundred times the difference between the replica-symmetric and one-step values, so enumeration cannot see replica symmetry breaking in the free entropy.

The script \texttt{01-freeentropy-curieweiss.py} computes the Curie--Weiss free entropy exactly up to $N=10^6$ in seconds and checks that every value lies within the bounds of \cref{thm:freeentropy:curieweiss} (\cref{fig:freeentropy:sk}(c)). The saddle-point corrections describe the data: at $\beta=0.5$ the values $0.326$, $0.344$, $0.3466$ at $N=10$, $10^2$, $10^4$ approach $-\frac12\log\frac12=0.3466$, and at $\beta=1.5$, where $m^*=0.8586$, the values $1.022$, $0.957$, $0.9440$ approach $0.9438$, which includes the $\log2$ of the two saddles. At $\beta=1$ the Gaussian correction is infinite and the data grow like $\frac14\log N$: the maximum is quartic, the magnetization lives on the scale $N^{-1/4}$ \parencite{EllisNewman1978}, and $Z_N/2^N\sim cN^{1/4}$ with $c=1.3460$ (\cref{ex:freeentropy:critical}), which matches the data to $0.0002$ at $N=10^6$.\sidenote{The script also checks \cref{cor:freeentropy:magnetization}: at $\beta=1.5$, $h=0.2$, $\gibbs m_N=0.7958$, $0.9143$, $0.9178$, $0.91814$ at $N=10$, $10^2$, $10^3$, $10^4$, against $m^*=0.91817$.}

\begin{table}[t]
\begin{fullwidth}
\centering\small
\setlength{\tabcolsep}{4.5pt}
\begin{tabular}{@{}cccccccc@{}}
\toprule
& \multicolumn{3}{c}{$\freeentropymmse(\gamma)$} & \multicolumn{2}{c}{Bayes posterior} & \multicolumn{2}{c}{posterior at SNR $2\gamma$}\\
\cmidrule(lr){2-4}\cmidrule(lr){5-6}\cmidrule(l){7-8}
$\gamma$ & exact & Monte Carlo & via I-MMSE & $\E[X\gibbs x]$ & $\E\gibbs x^2$ & $\E[X\gibbs x']$ & $\E\gibbs x'^2$\\
\midrule
$0.25$ & $0.7959$ & $0.7946(45)$ & $0.7942(88)$ & $0.2048$ & $0.2042$ & $0.2531$ & $0.3131$\\
$0.5$ & $0.6499$ & $0.6538(50)$ & $0.6527(85)$ & $0.3479$ & $0.3496$ & $0.4052$ & $0.4776$\\
$1$ & $0.4496$ & $0.4525(50)$ & $0.4562(76)$ & $0.5476$ & $0.5478$ & $0.6012$ & $0.6647$\\
$2$ & $0.2310$ & $0.2347(41)$ & $0.2329(61)$ & $0.7671$ & $0.7689$ & $0.8010$ & $0.8424$\\
$4$ & $0.0686$ & $0.0696(24)$ & $0.0726(40)$ & $0.9303$ & $0.9301$ & $0.9420$ & $0.9552$\\
\bottomrule
\end{tabular}
\end{fullwidth}
\caption{The scalar channel with the Rademacher prior. Exact: $1-\E\tanh(\gamma+\sqrt\gamma Z)$ by $161$-node Gauss--Hermite quadrature. Monte Carlo: mean squared error of $\tanh(\sqrt\gamma Y)$ over $10^5$ independent draws per $\gamma$ (seed $1000\gamma$), with the $95\%$ half-width in units of the last digits. Via I-MMSE: $2(\frac12-\psi'(\gamma))$, with $\psi'$ a centered finite difference (step $10^{-3}\gamma$) of the log-evidence on the same draws. Last four columns: the two sides of $\E[X\gibbs x]=\E\gibbs x^2$ ($95\%$ half-widths at most $0.003$) for the Bayes posterior and for a statistician who assumes the SNR $2\gamma$. Script: \texttt{figures/scripts/01-freeentropy-channel.py}.}\label{tab:freeentropy:channel}
\end{table}

The script \texttt{01-freeentropy-channel.py} checks the scalar channel in ten seconds (\cref{tab:freeentropy:channel}).\sidenote{By quadrature, $I'$ and $\frac12\freeentropymmse$ agree to ten digits at $\gamma=0.5$ and $2$, and so do $\E\tanh$ and $\E\tanh^2$.} Every Monte Carlo interval for the mean squared error contains the exact value, and so does every interval for the mmse computed from the derivative of the log-evidence alone: two measurements, of the evidence and of the error, tied by \cref{thm:freeentropy:immse}. The same holds for the Gaussian prior (twenty intervals in all). The Nishimori identity holds for the Bayes posterior within the Monte Carlo error and fails by a wide margin for the mismatched posterior, whose error also exceeds the mmse ($0.462$ against $0.450$ at $\gamma=1$): the identity is a property of the true posterior. No measurement contradicts a theorem of the chapter.

\section{The status board}\label{sec:freeentropy:status}

The status words are \emph{proved} (a theorem in a refereed publication, or classical), \emph{proved under} extra hypotheses that the board names, \emph{preprint} (proved only in an unrefereed preprint), \emph{conjectured} (with the evidence), \emph{heuristic} (an argument that claims no exactness), \emph{known false} (with the counterexample), and \emph{open}. As of October 2026 this board is nearly all proved; the one step that is not a computation, the continuation in the number of replicas, is not, the annealed answer fails for $\beta>1$, where the second moment grows exponentially, and the large deviations at low temperature rest on preprints.

\begin{statusboard}
\statusrow{Self-averaging of $\frac1N\log Z_N$ (Gaussian disorder)}{\proved}{Gaussian concentration \parencite{TsirelsonIbragimovSudakov1976,Borell1975}; \cref{thm:freeentropy:concentration}}
\statusrow{Curie--Weiss free entropy; $m=\tanh(\beta m+h)$}{\proved}{classical \parencite{Ellis1985}; \cref{thm:freeentropy:curieweiss}, error $\log(N+1)/N$}
\statusrow{Existence of $\lim\phi_N$ for SK}{\proved}{\textcite{GuerraToninelli2002}; \cref{thm:sk:guerratoninelli}}
\statusrow{Quenched $=$ annealed for SK, $h=0$, $\beta<1$}{\proved}{\textcite{AizenmanLebowitzRuelle1987}; \cref{ch:sk}}
\statusrow{Quenched $=$ annealed for SK, $h=0$, $\beta>1$}{\knownfalse}{Guerra's replica-symmetric bound is strictly smaller \parencite{Guerra2003}; \cref{ch:sk}}
\statusrow{Replica moments: the limits at integer $n$; the limit $n\to0$, uniformly in $N$}{\proved}{Laplace's method \parencite{Varadhan1966}, \cref{prop:freeentropy:moments}; concentration, \cref{prop:freeentropy:smalln}}
\statusrow{Replica trick: integer moments continued to $n\to0$}{\heuristicstatus}{no unconditional justification: counterexamples \parencite{Tanaka2007}; Carleman fails; the REM's exact moments are non-analytic \parencite{OgureKabashima2009a}; proof in \cref{thm:rem:moments,cor:rem:replicas}}
\statusrow{Upper tail $e^{-Na\delta^{6/5}}$ of $\frac1N\log Z_N$, SK, $h=0$, $\beta>1$}{\provedpreprint}{bounds of this order, and $\Var\log Z_N$ between the orders $N^{4/15}$ and $N^{7/15}$: arXiv:2609.09103v1 \parencite{AronowLopatto2026a}, written, its authors state, with the assistance of large language models; $a$ as $\delta\to0$: arXiv:2610.06480v2 \parencite{Pei2026}; predicted by \textcite{ParisiRizzo2009}}
\statusrow{Nishimori identity (Bayes form)}{\proved}{Bayes' rule; \cref{lem:freeentropy:nishimori}; \textcite[Proposition~16]{LelargeMiolane2019}}
\statusrow{I-MMSE relation $I'=\frac12\freeentropymmse$}{\proved}{\textcite{GuoShamaiVerdu2005}, finite variance; \cref{thm:freeentropy:immse}}
\end{statusboard}

\section{Where the translation fails}\label{sec:freeentropy:fails}

\paragraph{What was assumed.} The concentration theorem needs Gaussian disorder entering the Hamiltonian linearly with bounded coefficients; $\pm1$ couplings give the same tail by bounded differences (\cref{ex:freeentropy:rademacher}), heavy tails need other tools, and a planted signal is covered given the truth, which centers $\log Z_N$ at its conditional mean; for \eqref{eq:freeentropy:z2post} the symmetry of \eqref{eq:freeentropy:split} makes that mean free of the truth. The Curie--Weiss theorem uses the complete-graph symmetry and $\beta\ge0$ (\cref{ex:freeentropy:antiferro} treats $\beta<0$). The Nishimori identity needs the exact posterior: the statistician must know the prior, the channel and its signal-to-noise ratio. The I-MMSE relation needs additive Gaussian noise, and our proof bounded support. All limits are taken with $\beta$, $h$, $\lambda$, $\gamma$ fixed.

\paragraph{Where the physics is wrong, or not yet right.} The replica trick is not a theorem even at finite $N$: the identity $\E\log Z=\lim_{n\to0}(\E Z^n-1)/n$ holds (\cref{prop:freeentropy:smalln}), but integer moments do not determine the moments near $0$ \parencite{Tanaka2007}. In the random energy model the limit of the moments is known at every real $n>0$ and, for $\beta>\beta_c$, is not analytic at $n=\beta_c/\beta$, so the integer formula continued to $n\to0$ is wrong: \textcite{OgureKabashima2009a} derived the exact moment function, and \cref{thm:rem:moments,cor:rem:replicas} prove it. Annealed computations fail at low temperature, in the REM above $\beta_c$ and in SK above $\beta=1$, visibly so at $N=12$ in \cref{fig:freeentropy:sk}(b); the saddle-point correction fails at a critical point (panel~(c)). And not every quantity self-averages: in the low-temperature phase of SK physicists predict that the overlap distribution of a single sample fluctuates from sample to sample \parencite{MezardParisiVirasoro1987} (\cref{ch:sk}).

\paragraph{Open problems.} The order of the fluctuations of the free entropy at low temperature is open. For SK at $h=0$, $\beta<1$, $\log Z_N-\log\E Z_N$ has a Gaussian limit of order one \parencite{AizenmanLebowitzRuelle1987}. For $\beta>1$, \cref{thm:freeentropy:concentration} gives $\Var\log Z_N=O(N)$, \textcite{Chatterjee2009} proved $o(N)$ in an arXiv preprint, and \textcite{AronowLopatto2026a}, in an unrefereed preprint (arXiv:2609.09103v1) that its authors state was written with the assistance of large language models, prove that it lies between the orders $N^{4/15}$ and $N^{7/15}$. Replica theory predicts that $\frac1N\log Z_N$ exceeds its typical value by $\delta$ with probability $e^{-Na\delta^{6/5}}$ \parencite{ParisiRizzo2009}; the same preprint proves bounds of this order, and \textcite{Pei2026} (arXiv:2610.06480v2, unrefereed) identifies $a$ in the limit $\delta\to0$. Matching $Na\delta^{6/5}$ to one suggests a variance of order $N^{1/3}$, inside the window; this matching is our heuristic. More broadly, \textcite{Tanaka2007} notes that the replica method would be justified wherever the moment problem for $Z_N$ is solved affirmatively; we know of no general class of disordered systems in which this is proved.

\begin{exercises}
\item \consolidation (The critical point.) At $\beta=1$, $h=0$, use $s(m)=\log2-\frac{m^2}2-\frac{m^4}{12}+O(m^6)$ and the local central limit theorem to show $Z_N/2^N\sim cN^{1/4}$ with $c=2\cdot12^{1/4}\Gamma(\frac54)/\sqrt{2\pi}=1.3460$; compare with \cref{fig:freeentropy:sk}(c).\label{ex:freeentropy:critical}

\item \consolidation (Scalar channels.) For the standard Gaussian prior compute $\cZ_\gamma$, $\gibbs x$, $\psi$, $I$ and $\freeentropymmse$, and check \cref{thm:freeentropy:immse,cor:freeentropy:nishimori} directly; for the Rademacher prior derive $\psi(\gamma)$ and both expressions for the mmse.\label{ex:freeentropy:scalar}

\item \consolidation (Heyde's lognormal family.) Let $p_0$ be the density of $e^G$, $G\sim\Normal(0,1)$, and $p_a(z)=p_0(z)(1+a\sin(2\pi\log z))$, $\abs a\le1$. Show that each $p_a$ has the moments $e^{k^2/2}$, $k=0,1,\dots$, and the mean logarithm $2\pi ae^{-2\pi^2}\approx1.68\times10^{-8}a$. What does the replica trick return? \hint{$\E[e^{kG}f(G)]=e^{k^2/2}\E f(G+k)$.}\label{ex:freeentropy:heyde}

\item \consolidation (A corner.) For the Curie--Weiss model at $h=0$, $\beta>1$, show that $\partial_h\phi(\beta,0\pm)=\pm m^*(\beta)$ and $\gibbs{\abs m}_N\to m^*(\beta)$. \hint{bound the sum over types with $\abs{\abs m-m^*}\ge\epsilon$ as in \cref{lem:freeentropy:types}.}\label{ex:freeentropy:corner}

\item \byhand (Three replicas.) Show that $\Lambda_3(\beta)\ge\max(3\log2,\log2+\frac32\beta^2)$, using the uniform law of the column and the law with mass $\frac12$ on each of $\pm(1,1,1)$; the second is larger exactly when $\beta>(\frac43\log2)^{1/2}=0.9614$. At $\beta=1.5$ compare the lower bound $\frac13(\frac34\beta^2+\log2+\frac32\beta^2)=1.9185$ on the limit with the exact $1.8167$ at $N=12$, using the error terms of \cref{prop:freeentropy:moments}.\label{ex:freeentropy:threereplicas}

\item \byhand (The random-field Curie--Weiss model.) Let $H_N=\frac\beta{2N}S^2+\sum_ih_i\sigma_i$ with $\beta>0$ and $h_i$ independent copies of a bounded $h$. (a) Rewrite \eqref{eq:freeentropy:hs} with $\sum_i\log2\cosh(\beta\mu+h_i)$ in place of $N\log2\cosh(\beta\mu+h)$, and show with the law of large numbers that, almost surely, $\frac1N\log Z_N$ tends to $\max_\mu[\E\log2\cosh(\beta\mu+h)-\frac\beta2\mu^2]$. (b) Compute $\E Z_N^n$ with one auxiliary field per replica, take the fields equal, and check that $n\to0$ gives (a).\label{ex:freeentropy:rfcw}

\item \byhand (Mismatched posteriors.) For the standard Gaussian prior, a statistician who assumes the SNR $\gamma'$ uses $\hat x=\sqrt{\gamma'}Y/(1+\gamma')$ and posterior variance $1/(1+\gamma')$. Show that, for $\gamma'>0$, $\E[X\hat x]=\E\hat x^2$ exactly when $\gamma'=\gamma$ or $\gamma\gamma'=1$, while the posterior second moment has mean $\E X^2$ only when $\gamma'=\gamma$.\label{ex:freeentropy:mismatch}

\item \translate (A saddle point that is a minimum.) For $\beta<0$ show that $\phi+\frac\beta{2N}\le\frac1N\log Z_N\le\phi+\frac{\log(N+1)}N$, with $\phi=\max_m\Phi(m)$ as in \cref{thm:freeentropy:curieweiss}, that the integral \eqref{eq:freeentropy:hs} diverges, and, by the minimax theorem, that $\phi=\min_\mu[\log2\cosh(\beta\mu+h)-\frac\beta2\mu^2]$. Compare with the replica saddle point for $n<1$.\label{ex:freeentropy:antiferro}

\item \translate\hard (Quenched equals annealed from two moments.) Suppose a Gaussian disordered system of size $N$ has $L^2\le cN$ and $\E Z_N^2\le C(\E Z_N)^2$ for all $N$. From the Paley--Zygmund inequality $\Prob(Z\ge\frac12\E Z)\ge(\E Z)^2/(4\E Z^2)$ and \cref{thm:freeentropy:concentration}, show $\E\log Z_N\ge\log\E Z_N-\log2-L\sqrt{2\log(4C)}$, and deduce from \cref{comp:freeentropy:secondmoment} that $\phi=\log2+\beta^2/4$ for SK at $h=0$, $\beta<1$. \hint{argue by contradiction on the event $\{Z_N\ge\frac12\E Z_N\}$.}\label{ex:freeentropy:secondmoment}

\item \translate (The vector I-MMSE relation.) For a prior on a bounded subset of $\R^N$ and $Y=\sqrt\gamma X+Z$ with $Z\sim\Normal(0,I_N)$, adapt the proof of \cref{thm:freeentropy:immse} to show $\frac{\dd}{\dd\gamma}I(X;Y)=\frac12\E\norm{X-\gibbs x}^2$, and write the derivative of the free entropy as an overlap.\label{ex:freeentropy:vector}

\item \measurement (Rademacher couplings.) (a) For SK with $g_{ij}$ uniform on $\{\pm1\}$, show by the bounded-differences inequality that \eqref{eq:freeentropy:sktail} still holds. (b) Modify \texttt{01-freeentropy-sk.py} accordingly: at $N=12$, $\beta=0.5$, the variance of $\log Z_N$ drops from about $0.016$ to about $0.002$, while $\phi_N$ barely moves. Explain the drop by the term $\frac{\beta^2}{2N}\sum_{i<j}(g_{ij}^2-1)$, present only for Gaussian couplings, whose variance is $\beta^4(N-1)/(4N)$.\label{ex:freeentropy:rademacher}

\item \research (The low-temperature exponent.) Explain why matching $Na\delta^{6/5}$ to one suggests $\Var\log Z_N\asymp N^{1/3}$ for SK at $h=0$, $\beta>1$. Estimate the exponent by enumeration and Monte Carlo, and explain why accessible sizes cannot separate $\frac4{15}$, $\frac13$ and $\frac7{15}$. Determining it is open.\label{ex:freeentropy:exponent}
\end{exercises}

\notesandrefs
\Textcite{EllisNewman1978} proved limit theorems for the Curie--Weiss magnetization, with a non-Gaussian law for $S/N^{3/4}$ at the critical point; \textcite{Ellis1985} gives the large-deviation treatment, and Varadhan's lemma is \textcite{Varadhan1966}. Gaussian concentration with the sharp constant is due to \textcite{TsirelsonIbragimovSudakov1976} and \textcite{Borell1975}; the appendix proves it with a weaker constant and states the sharp one (\cref{thm:toolkit:lipschitz}), and for log-partition functions the covariance identity \eqref{eq:freeentropy:covariance}, which goes back to \textcite{HoudrePerezAbreu1995}, gives the sharp constant directly, as in part~2 of \cref{thm:freeentropy:concentration}. The proofs of \cref{thm:freeentropy:curieweiss,thm:freeentropy:concentration,prop:freeentropy:moments,prop:freeentropy:smalln} are written for this book from standard ingredients; we claim no novelty for them.

Replicas entered spin-glass theory with \textcite{EdwardsAnderson1975} and the solvable model of \textcite{SherringtonKirkpatrick1975}; their failure at low temperature led to Parisi's broken replica symmetry \parencite{Parisi1979b,MezardParisiVirasoro1987}. The moments at real $n>0$ describe large deviations of the free energy \parencite{ParisiRizzo2009} and for $0<n<1$ were related by \textcite{Talagrand2007} to a Parisi problem with an atom of mass at least $n$ at $0$ (as \textcite{Pei2026} reports); \textcite{Tanaka2007} sets out the objections to the continuation, with the lognormal family of \textcite{Heyde1963} among his examples. The limit for SK exists by \textcite{GuerraToninelli2002}, and Guerra's bounds \parencite{Guerra2003} open the rigorous theory of \cref{ch:sk}.

The Nishimori line goes back to \textcite{Nishimori1981} and its Bayesian reading to \textcite{Iba1999}, both developed in \textcite{Nishimori2001}. The I-MMSE relation is \textcite{GuoShamaiVerdu2005}. \Textcite{ZdeborovaKrzakala2016} build the statistical physics of inference on these identities, and \textcite{MontanariSen2024} explain replica computations to non-physicists.

\paragraph{Translating the literature.} Physicists write the Gibbs weight $e^{-\beta\mathcal H}$, with $T=1/\beta$ and free energy per spin $f=-\phi/\beta$; \textcite{SherringtonKirkpatrick1975} take couplings of mean $J_0/N$ and variance $J^2/N$, so our $\beta$ is their $\beta J$ and their $\beta J_0$ is our Curie--Weiss coupling. \Textcite{Talagrand2011a} writes $p_N(\beta,h)$ for our $\phi_N$, $R_{1,2}$ for our $R_{12}$, and $\nu(f)$ for $\E\gibbs f$; \textcite{Panchenko2013b} calls $\frac1N\E\log Z_N$ the free energy, without the sign and the factor $1/\beta$ of physics, and writes $\xi(q)=\beta^2q^2$, so his $\beta$ is our $\beta/\sqrt2$. \Textcite{GuoShamaiVerdu2005} write $\mathrm{snr}$ for our $\gamma$; \textcite{LelargeMiolane2019} write $Y=\sqrt{\lambda/n}\,XX^\top+Z$, so their $\lambda$ is our $\lambda^2$.

\printchapterbib
\end{refsection}
